% =====================================================================
%  TomaLayer: the block, and every stage it is instantiated as.
%  Top row  = the shared backbone stages (TomaNet-C and TomaNet-D).
%  Bottom row = the neck fusion points (TomaNet-D only).
%  Tile height tracks feature-map size, tile width tracks channels, so
%  the funnel shape matches the real models. Arguments under each tile
%  are the literal YAML arguments from configs/models/tomanet-*.yaml.
% =====================================================================
\documentclass[border=14pt, tikz]{standalone}
\usepackage{import}
\subimport{./layers/}{init}
\usetikzlibrary{positioning,3d,calc,backgrounds}
% =====================================================================
%  tomanet_palette.tex -- one colour per architectural role, shared by
%  fig_tomanet_c.tex and fig_tomanet_d.tex so a shared module (the stem
%  and backbone) keeps the exact same colour in both figures. That
%  visual identity IS the point: TomaNet-C and TomaNet-D use a
%  byte-identical trunk (layers 0-8) and only the colours after the
%  trunk change.
% =====================================================================
\definecolor{tStem}   {RGB}{243,203,150}  % stem convs             sand
\definecolor{tStemP}  {RGB}{253,240,222}  % stem panel
\definecolor{tTrunk}  {RGB}{150,196,140}  % TomaLayer stages        leaf green
\definecolor{tTrunkP} {RGB}{232,244,227}  % trunk panel
\definecolor{tBlkA}   {RGB}{120,175,190}  % RepPConv / spatial      teal
\definecolor{tBlkB}   {RGB}{150,150,200}  % CoordAtt                periwinkle
\definecolor{tNeck}   {RGB}{ 90,110,175}  % TomaLayer, neck (e=1.0) deep slate blue
\definecolor{tNeckP}  {RGB}{231,236,247}  % neck panel
\definecolor{tStruct} {RGB}{200,200,200}  % structural ops: upsample / concat -- not a learned layer, so deliberately neutral grey
\definecolor{tHeadC}  {RGB}{210,110,90}   % classify head           tomato red
\definecolor{tHeadCP} {RGB}{250,231,225}  % classify head panel
\definecolor{tHeadD}  {RGB}{210,110,90}   % detect heads            tomato red (same family: both are "the answer")
\definecolor{tHeadDP} {RGB}{250,231,225}  % detect head panel
\definecolor{tLine}   {RGB}{ 90,100,95}   % connector
\definecolor{tRule}   {RGB}{200,206,200}  % separators

% ---- PlotNeuralNet 3D-tile band colours (darker shade of each block's
% own hue, used for RightBandedBox's second face -- e.g. "the BN+act
% half" or "the attention half" of a two-part operation) -------------
\colorlet{tStemBand}  {tStem!60!black}
\colorlet{tTrunkBand} {tTrunk!60!black}
\colorlet{tNeckBand}  {tNeck!60!black}
\colorlet{tHeadCBand} {tHeadC!60!black}
\colorlet{tHeadDBand} {tHeadD!60!black}
\colorlet{tBlkABand}  {tBlkA!60!black}
\colorlet{tBlkBBand}  {tBlkB!60!black}
\definecolor{tBlkSum} {RGB}{190,140,150}  % residual / concat ball    dusty rose


\begin{document}
\begin{tikzpicture}[font=\Large]

\tikzstyle{connection}=[ultra thick,draw=\edgecolor,opacity=0.8]
\tikzstyle{arr}=[connection,-{Stealth[length=3.6mm,width=2.8mm]}]
\tikzstyle{arrsoft}=[connection,densely dashed,opacity=0.55,
                     -{Stealth[length=3.0mm,width=2.4mm]}]
\tikzstyle{ttl}=[font=\LARGE\bfseries,text centered]
\tikzstyle{sub}=[font=\Large,text centered]
\tikzstyle{grouptitle}=[font=\Large\bfseries,anchor=west]
\tikzstyle{args}=[font=\large,text centered,anchor=north,text width=4.4cm]
\tikzstyle{note}=[font=\large,text centered,anchor=north,text width=5.2cm]

% ----------------------------------------------------------------
% background panels for the two roles
% ----------------------------------------------------------------
\begin{scope}[on background layer]
  \fill[tTrunk!20,rounded corners=6pt] (16.4,10.2) rectangle (40.4,1.5);
  \fill[tNeck!16,rounded corners=6pt]  (16.4,-4.3) rectangle (40.4,-11.7);
\end{scope}

% ----------------------------------------------------------------
% title
% ----------------------------------------------------------------
\node[ttl] at (20,14.4) {TomaLayer: One Block, Eight Stages};
\node[sub,text width=30cm] at (20,13.2)
     {\textbf{n} TomaBlocks stacked $=$ one TomaLayer. Same code every
      time -- only the arguments below change.};
\node[sub,font=\large,text=black!65,text width=30cm] at (20,12.2)
     {Values are the YAML arguments, before the n/s/m multipliers scale
      channels and n (at scale n, stage 2 builds 32 channels, 1 block).};

% ----------------------------------------------------------------
% the block and the layer
% ----------------------------------------------------------------
\pic at (0,0,0)
     {Box={name=blk,caption={TomaBlock},xlabel={{"",""}},%
       fill=tBlkA,opacity=0.9,height=22,width=18,depth=7}};
\node[note] at (1.8,-3.9) {the unit: space, channels, attention};

\pic[shift={(4.0,0,0)}] at (blk-east)
     {Box={name=lay,caption={TomaLayer},xlabel={{"",""}},%
       fill=tBlkB,opacity=0.9,height=22,width=18,depth=7}};
\node[note] at (9.4,-3.9) {n of them, back to back};

\draw[arr] (blk-neareast) -- node[above,font=\large]{$\times$\,n} (lay-nearwest);

% ----------------------------------------------------------------
% BACKBONE ROW -- 4 stages, shared by both models
% ----------------------------------------------------------------
\node[grouptitle,text=black!75] at (16.9,10.9)
     {Backbone -- the 4 shared stages (TomaNet-C \emph{and} TomaNet-D)};

\pic at (18,7,0) {Box={name=s2,caption={stage 2},xlabel={{"128",""}},%
       fill=tTrunk,opacity=0.9,height=26,width=13,depth=7}};
\pic at (24,7,0) {Box={name=s3,caption={stage 3},xlabel={{"256",""}},%
       fill=tTrunk,opacity=0.9,height=20,width=15,depth=7}};
\pic at (30,7,0) {Box={name=s4,caption={stage 4},xlabel={{"512",""}},%
       fill=tTrunk,opacity=0.9,height=14,width=17,depth=7}};
\pic at (36,7,0) {Box={name=s5,caption={stage 5},xlabel={{"1024",""}},%
       fill=tTrunk,opacity=0.9,height=9,width=19,depth=7}};

\node[args] at (19.3,2.9) {k=3, n=2\\e=2.0};
\node[args] at (25.5,2.9) {k=3, n=4\\e=2.0};
\node[args] at (31.7,2.9) {k=5, n=4\\e=2.0};
\node[args] at (37.9,2.9) {k=5, n=2\\e=2.0};

\draw[arrsoft] (s2-neareast) -- (s3-nearwest);
\draw[arrsoft] (s3-neareast) -- (s4-nearwest);
\draw[arrsoft] (s4-neareast) -- (s5-nearwest);

% ----------------------------------------------------------------
% NECK ROW -- 4 fusion points, detector only
% ----------------------------------------------------------------
\node[grouptitle,text=black!75] at (16.9,-3.6)
     {Neck -- 4 more fusion points (TomaNet-D only, lighter e)};

\pic at (18,-7,0) {Box={name=k1,caption={P4 fuse},xlabel={{"512",""}},%
       fill=tNeck,opacity=0.9,height=14,width=17,depth=7}};
\pic at (24,-7,0) {Box={name=k2,caption={P3 small},xlabel={{"256",""}},%
       fill=tNeck,opacity=0.9,height=20,width=15,depth=7}};
\pic at (30,-7,0) {Box={name=k3,caption={P4 medium},xlabel={{"512",""}},%
       fill=tNeck,opacity=0.9,height=14,width=17,depth=7}};
\pic at (36,-7,0) {Box={name=k4,caption={P5 large},xlabel={{"1024",""}},%
       fill=tNeck,opacity=0.9,height=9,width=19,depth=7}};

\node[args] at (19.7,-10.5) {k=3, n=2\\e=1.0};
\node[args] at (25.5,-10.5) {k=3, n=2\\e=1.0};
\node[args] at (31.7,-10.5) {k=3, n=2\\e=1.0};
\node[args] at (37.9,-10.5) {k=5, n=2\\e=1.0};

\draw[arrsoft] (k1-neareast) -- (k2-nearwest);
\draw[arrsoft] (k2-neareast) -- (k3-nearwest);
\draw[arrsoft] (k3-neareast) -- (k4-nearwest);

% ----------------------------------------------------------------
% fan-out: one layer definition, two roles
% ----------------------------------------------------------------
\coordinate (fork) at ($(lay-neareast)+(1.6,0)$);
\draw[connection] (lay-neareast) -- (fork);
\draw[arr] (fork) |- (s2-nearwest);
\draw[arr] (fork) |- (k1-nearwest);

\node[sub,font=\large,text width=5cm] at (13.9,1.9) {built as};

\end{tikzpicture}
\end{document}
